\documentclass[11pt]{article}
\usepackage[a4paper,margin=27mm]{geometry}
\usepackage[T1]{fontenc}
\usepackage{lmodern,amsmath,amssymb,amsthm,mathtools,booktabs,longtable,array}
\usepackage{microtype}
\usepackage[unicode,colorlinks=true,linkcolor=blue!45!black,citecolor=blue!45!black,urlcolor=blue!45!black]{hyperref}
\usepackage{xcolor}
\newtheorem{theorem}{Theorem}[section]
\newtheorem{lemma}[theorem]{Lemma}
\newtheorem{proposition}[theorem]{Proposition}
\newtheorem{remark}[theorem]{Remark}
\newcommand{\eps}{\varepsilon}
\newcommand{\cm}{\operatorname{cm}}
\newcommand{\OPT}{\operatorname{OPT}_{\rm RC}}
\newcommand{\pref}{\operatorname{pref}}
\newcommand{\suf}{\operatorname{suf}}
\newcommand{\outdeg}{\operatorname{outdeg}}
\newcommand{\indeg}{\operatorname{indeg}}
\newcommand{\ZZ}{\mathbb Z}
\newcommand{\src}[1]{\cite[#1]{ordinary}}
\setlength{\emergencystretch}{2em}
\title{A source-relative factor-two approximation\protect\\for reverse-complement superstrings}
\author{Myeongjin Kang}
\date{Manuscript r1, October 2026}
\begin{document}
\maketitle
\begin{abstract}
Assuming the count, periodic-layer, and connection lemmas of the pinned
ordinary shortest-common-superstring construction of OpenAI, we give a
deterministic polynomial-bit-time factor-two algorithm for superstrings
with reverse complements. The alphabet and its involution are explicitly
encoded and may vary with the input; fixed letters, reverse-complement
palindromes, empty requirements, duplicates, and containment are included.
The extension uses one matched construction: a phase-free $2k+1$ count
rule for self-reversed periodic texts, adjacent sorted profiles sharing
both reflection axes, and an odd core of $2k-1$ layers in a $2k$-layer
budget unit. Established connection and symmetric traversal operations
complete the algorithm. The dependence is on the stated source's stronger
structural lemmas, not merely on black-box existence of an ordinary
factor-two algorithm. Sixteen finite and scalar Lean support lemmas cover
specified interfaces only; the full algorithm is not formalized.
\end{abstract}
\section{Statement, dependence, and antecedents}\label{sec:scope}
Let $\Sigma$ be a finite alphabet with an involution $\cm:\Sigma\to\Sigma$.
For a word $s=s_1\cdots s_h$, put
$\rho(s)=\cm(s_h)\cdots\cm(s_1)$ and $\rho(\eps)=\eps$.
A word is feasible for a finite listed family $\mathcal S$ if, for every
$s\in\mathcal S$, it contains $s$ or $\rho(s)$ as a contiguous substring.
Write $\OPT(\mathcal S)$ for its minimum symbol length. Alphabet labels,
the involution table, word symbols, and list boundaries are explicit;
$N$ denotes their total encoded bit length.

\begin{theorem}[Source-relative approximation]\label{thm:main}
Assume the structural results of \cite{ordinary} at revision
\textnormal{\texttt{fd4aeeb2ee4fc729c18d98444fed42fd0529eeeb}}, with the contracts
specified below. There is a deterministic algorithm, polynomial in $N$,
which for every such $(\Sigma,\cm,\mathcal S)$ returns a feasible word $T$
with $|T|\le2\OPT(\mathcal S)$. The assertion includes arbitrary fixed
letters and words, empty words and the empty list, duplicates, and substring
redundancy. Neither the alphabet nor the word lengths are bounded constants.
\end{theorem}

The source's truths are assumptions of this article. We do not audit them
or derive an independent proof of the ordinary superstring theorem.
In particular, an unspecified ordinary factor-two algorithm would not
supply the contracts used here. Yamano--Shibuya's reduction in
\cite[Theorem 2]{yswabi} also explains why an unrestricted reverse-complement
result cannot be advertised as bypassing the ordinary problem.

\subsection{The precise source interface}\label{sec:contracts}
The following labels identify results in the pinned source, not labels
of this article. Its section files are under \texttt{build/sections/}.
\begin{enumerate}
\item \path{02-counts.tex}: the blocking functional, extension and
ordinary period rules, and the occurrence and balance proofs
\path{lem:counts}, \path{lem:base}. Their use requires a nonnegative
count vector on a substring-closed set satisfying those rules. It does
not require choosing the smallest ordinary count vector.
\item \path{03-layers.tex}: \path{lem:recorded}, \path{lem:layers},
\path{lem:matching-window}, \path{lem:agreement}, and
\path{lem:max-rotation}. These provide recorded counts, sorted primitive
passages, actual matching windows, and the two periodic comparison bounds.
\item \path{04-connections.tex}: \path{lem:ordered-windows} and
\path{lem:band}; the \emph{proof} of \path{lem:collective}; and
\path{lem:host} after its hypotheses are established on a single host
path. We state the relevant contracts in Section~\ref{sec:operations}.
The source collective statement has $k$ layers and $K=kp$; our odd-core
use has $2k-1$ layers and $K=kp$, and uses the proof's geometric inequalities.
\item \path{05-processing.tex}: its ordinary threshold, partition,
and individual-link/request construction, applied to nonself text orbits.
\item \path{06-cycles.tex}: the proof of \path{lem:cycle-contact}
and the free-period strategy of \path{lem:free-period} and
\path{thm:connection}. We adapt the geometric proof to odd cores and
justify the two-port quotient separately; the original one-sided
connection theorem is not assumed equivariant.
\item \path{07-algorithm.tex}: finite record search, explicit coordinate
representation, and the suffix invariant for spelling the output.
\end{enumerate}
All these hypotheses and their symmetric applications are made explicit
below. Source-paid arguments are recalled in our notation where needed to
make the adaptation checkable.

\subsection{Related operations and the credit boundary}\label{sec:antecedents}
The arbitrary-involution formulation, formal copies of fixed words,
paired cycle selection, both-orientation periodic comparisons, and an
$8/3$ approximation occur in Yamano--Shibuya \cite{yswabi}, especially
Section 2, Lemmas 7, 9, 12, and Theorems 11--13. Their CPM paper
\cite{yscpm} proves $3.75$ and $2.875$ guarantees and, in Lemma 18,
compares an optimum with its entire reverse complement. Thus the generic
use of $T$ and $\rho(T)$ is an antecedent, not a contribution here.
Fici et al.\ \cite{fici} already use distinct palindrome vertex copies
and paired reverse paths in Sections 2--3, and reversal surgery in
Lemma 6. Their Theorem 7 is a one-half \emph{compression} guarantee;
it is not a factor-two length guarantee.

Jiang, Li and Du studied superstrings with flipping, as described by
Fici et al.\ and Yamano--Shibuya \cite{jld,fici,yscpm,yswabi}.
The 1992 full text was unavailable in our bounded searches, so our
comparison uses those later accounts. Fici et al.\ report a $4$ length
bound for a greedy-like reversal algorithm. The CPM 2026 account
distinguishes the $4$-approximate MGREEDY extension from a suggested
$3$-approximate TGREEDY extension; WABI 2026 summarizes the earlier bound
as $3$. We do not assign an independent theorem statement to the unread
1992 body.

Goldberg--Karzanov's symmetric decomposition and parity results
\cite[Theorem 3.1 and Corollary 3.2]{gk} precede the pairing operations
below. Their Section 5 already exchanges complementary halves of two
self-symmetric cycles meeting at two mate vertices. We credit that splice;
our shared-axis lemma supplies the meeting hypothesis for possibly
different ordered staircase profiles. Their setting has no fixed vertices,
whereas our hierarchical graph has fixed vertices, including $\eps$;
we therefore give the needed parity and traversal proof explicitly.
Bidirected incidence walks and balanced connected Euler criteria also
appear in Medvedev et al.\ \cite[Section 2.2, Lemma 1]{medvedev}
and Kundeti et al.\ \cite[Section 6]{kundeti}. Their spectrum-constrained
cyclic assembly results are not imported as exact algorithms for our
arbitrary-length linear requirements.

Extremal and critical rotations are established tools in
Breslauer--Jiang--Jiang \cite{bjjjournal,bjjreport} and Mucha \cite{mucha}.
We cite the 1997 journal article as the archival reference, but our
operational comparison uses the 1996 BRICS report RS-96-21; the journal
body was not read, and report lemma numbers are not transferred to it.
The report's concluding remarks credit Crochemore with a Lyndon
minimum/maximum-rotation argument. The hierarchical representation and
the separation between cycle-cover bounds and connectivity also precede
this work; see the Collapsing framework \cite[Sections 2.2 and 7.2]{collapsing}.

The contribution claimed here is exactly one matched package:
the phase-free self-text count rule in Lemma~\ref{lem:domination}, the transfer
to adjacent sorted profiles sharing both axes in Lemma~\ref{lem:axes}, and
the $2k-1$ odd core inside a $2k$-layer unit in Lemma~\ref{lem:odd-core}.
No separate novelty is claimed for phase parity, tagged counts,
maxima or even rounding, decomposition, splicing, anchors, batch flips,
request reclassification, rotations, cycle accounting, half-Euler
extraction, or complexity bookkeeping.

\subsection{Why the matched package is needed}\label{sec:motivation}
For identity complement on $\{b,c\}$, take
$\mathcal S=\{bbcbbc,bcbbcb,cbbcbb\}$. The ordinary rules symmetrized
by orbit maxima and even rounding give one-letter counts $4,2$, hence
$W_{\rm sym}=6$. But $\OPT=7$: a length-six word cannot contain the two
distinct length-six orientation orbits, and $bbcbbcb$ satisfies both.
This refutes only a connection theorem that would retain those counts
and yield output length at most the same $W_{\rm sym}$. It does not
refute factor two or all equivariant methods. The stronger rule below
gives $W=12$ on this instance.

A different diagnostic, $\{aaa,aab,abcaa,accc\}$ with $a\leftrightarrow b$
and $c$ fixed, has optimum $10$. Minimum constrained-cycle-cover
orientation ties can lead to ordinary forced-count certificates $11$
or $9$. This shows that arbitrary minimum tie-breaking does not ensure
the desired $W$ certificate; it does not exclude all orientation-selection
methods. These finite certificates motivate the construction and are not
used to prove Theorem~\ref{thm:main}.

Boosting only fixed words is also insufficient for the intended proof:
in $(abcacb)^{\ZZ}$, the layer $z(x)=x+8$ has fixed-axis windows at starts
$0,3$, whereas its unique greatest forward rotation starts at $4$.
The two windows there are $cbabcac$ and $cbabcacb$, neither fixed.
A count rule confined to fixed-axis starts would not justify the
maximal-rotation cycle argument. The rule below holds at every phase
of a self-reversed text.

\section{Symmetric counts and their domination}\label{sec:counts}
Remove empty requirements, merge $\rho$-orbits, and discard an orbit
contained in a different larger orbit. If $s$ is contained in $t$, then
$\rho(s)$ is contained in $\rho(t)$, so either orientation of the larger
word satisfies an orientation of the smaller. This preprocessing preserves
feasibility and the optimum. If nothing remains, return $\eps$.
Otherwise let $R$ include both orientations of the surviving requirements,
and let $V$ consist of all their substrings, including $\eps$.
It is substring-closed and $\rho$-closed; required words are maximal in $V$.
Let $L$ be the total symbol length of one representative per surviving
orbit, so $L\le N$ and $|V|=O(L^2)$.

\subsection{Blocking and the recursion}\label{sec:blocking}
A periodic text is a map $A:\ZZ\to\Sigma$ with $A(i+p)=A(i)$ for
some positive $p$, not necessarily least. For $s=A[0:h)$ and $r\in V$,
let $b_{A,s}(r)$ count the integers $j$ such that
\begin{gather}
1\le j,\quad j+h\le |r|-1,\quad r[j:j+h)=s,\notag\\
r(i)=A(i-j)\quad(1\le i\le|r|-2),\qquad
r(0)\ne A(-j),\quad r(|r|-1)\ne A(|r|-1-j).\label{eq:block}
\end{gather}
Put $B_A(s;m)=\sum_{r\in V\setminus\{\eps\}}b_{A,s}(r)m(r)$.
Only $|r|\ge |s|+2$ contributes. For an occurrence aligned with $A$
in a finite word, the nearest mismatches on both sides, if present,
determine a unique bracket. Call it blocked when that bracket belongs
to $V$; absence of either mismatch means unblocked. Consequently
$B_A(s;N_T)$ counts blocked occurrences exactly. This is the blocking
contract of \src{Section 2}; its coefficients count placements,
not just distinct bracket words.

Define $m(s)$ in decreasing length, simultaneously on each $\rho$-orbit.
Missing words have count zero. The lower bounds are
\begin{equation}\label{eq:extension}
m(s)\ge\sum_c m(cs),\qquad m(s)\ge\sum_c m(sc),
\end{equation}
with required lower bound $1$ for nonfixed words and $2$ for fixed words.
For $v=A[0:h)$, $w=A[0:h-kp)$ nonempty, $k\ge1$, impose
\begin{equation}\label{eq:new-rule}
\begin{aligned}
&m(v)>B_A(v;m)\quad\Longrightarrow\quad m(w)\ge B_A(w;m)+g(A,k),\\
&g(A,k)=\begin{cases}2k+1,&A\text{ self-reversed},\\ k+1,&\text{otherwise}.
\end{cases}
\end{aligned}
\end{equation}
Here self-reversed means that for some integer $c$,
$\cm(A(j))=A(c-1-j)$ for every $j$; equivalently its reverse complement
is a translate of itself. Take the maximum of all bounds for both words
of an orbit and assign it to both. At a fixed word round upward to even.
Every bound refers only to longer words, so this defines an explicit
recursion. Enumerate every $v\in V$, $1\le p<|v|$, and $k$ with
$kp<|v|$, retaining a rule precisely when repetition of $v[0:p)$ agrees
with $v$. This includes all nonleast periods and all phases represented
by such words. No fixed-word-only $2k+2$ variant is used.

\begin{lemma}[Tagged-text domination]\label{lem:domination}
For every feasible $T$ and every nonempty $s\in V$,
\[
m(s)\le N(s):=N_T(s)+N_{\rho(T)}(s).
\]
Thus $W:=\sum_{a\in\Sigma}m(a)\le2\OPT(\mathcal S)$, and $m(s)\le2L$.
\end{lemma}
\begin{proof}
Treat $T$ and $\rho(T)$ as distinct \emph{tagged components}, even if
their word labels coincide. Then $N$ is symmetric and even at every
fixed word. It satisfies the required and extension bounds. Descending
induction for the ordinary $k+1$ rule is the proof of
\src{\texttt{lem:counts}} on this disjoint collection: an unblocked
$v$ gives $k+1$ period-spaced unblocked prefixes, or a positive
longer-bracket deficit contributes that many placements. The blocking
functional is linear over the components.

We prove the additional $2k+1$ bound. Translate so that $v=A[0:h)$,
$|w|=h-kp$. In a self-reversed text, $\rho(v)$ occupies $A[c-h:c)$.
The possible starts of a copy of $w$ inside this interval form an
integer interval of $kp+1$ integers. At least $k$ of them have residue
$0$ modulo $p$, and hence spell $w$ in the required $A$ alignment.
The original $v$ contains $k+1$ such copies, at $0,p,\ldots,kp$.

If $N(v)>B_A(v;N)$, choose an unblocked $v$ in one tagged component.
All $k+1$ internal copies are unblocked: a mismatch bracket of an
internal copy must enclose the entire matching $v$ interval. Reflect
this occurrence into the other component. Each of the indicated $k$
aligned copies in $\rho(v)$ is likewise unblocked. A bracket blocking
one would enclose $\rho(v)$; its reflection is in $V$ and would block
the original $v$. End mismatches remain mismatches because $\cm$ is a
bijection. Tagging makes the two sets of occurrences disjoint. Therefore
\[
N(w)-B_A(w;N)\ge(k+1)+k=2k+1.
\]

Otherwise $N(v)=B_A(v;N)\ge m(v)>B_A(v;m)$. Write
$\Delta(r)=N(r)-m(r)$ for longer words. Induction makes these
nonnegative. Some bracket word $r$ has $b_{A,v}(r)>0$ and
$\Delta(r)>0$. Its chosen placement contains $k+1$ counted placements
of $w$. The reflected bracket in $\rho(r)$ contains at least $k$
of the aligned shifted-phase placements just described; both endpoints
remain mismatches, so they are counted by the same functional for $w$.
If $r\ne\rho(r)$, symmetry gives
$\Delta(\rho(r))=\Delta(r)\ge1$, and these two bracket terms contribute
at least $2k+1$ to $B_A(w;N)-B_A(w;m)$. If $r=\rho(r)$, the positive
integer $\Delta(r)$ is even, hence at least $2$; the original $k+1$
placements alone give $2(k+1)\ge2k+1$. There is no double counting of
a fixed bracket. All other terms are nonnegative. This proves the new
bound in both cases.

Orbit maxima stay below the symmetric $N$, and even rounding stays below
its even values at fixed words. Induction follows. Summing letter
counts gives $W\le2|T|$. Concatenating one surviving representative per
orbit supplies a feasible word of length $L$, proving $m(s)\le2L$.
\end{proof}

\subsection{The balanced graph and fixed vertices}\label{sec:base}
For each nonempty $s$, the hierarchical graph has an up edge
$\pref(s)\to s$ and a down edge $s\to\suf(s)$; prefix and suffix
here delete the last and first letter respectively. Up cost is one,
down cost zero. Give these edges multiplicities
\begin{equation}\label{eq:multiplicities}
u(s)=m(s)-\sum_c m(cs),\qquad d(s)=m(s)-\sum_c m(sc).
\end{equation}
By the telescoping proof of \src{\texttt{lem:base}}, this graph $G$ is
balanced, has $W$ up and $W$ down copies, and uses every required vertex.
Indeed both degrees at nonempty $s$ equal
$m(s)-\sum_{c,d}m(csd)$, and the total multiplicities telescope to $W$.

Symmetry gives $u(s)=d(\rho(s))$. Pair copies to define an
edge-reversing involution $\tau$, mapping the up edge at $s$ to the
down edge at $\rho(s)$. There are no fixed edges. At a nonempty fixed
$s$, the displayed degree is even: nonfixed extension words pair and
fixed extension words have even counts. At the empty vertex,
\begin{equation}\label{eq:empty-degree}
\outdeg_G(\eps)=\sum_a u(a)
=\sum_{|s|=1}m(s)-\sum_{|s|=2}m(s)
\end{equation}
is also even. It is not the total up cost $W$. Adding a closed walk
$C$ together with its reversed image $\tau C$ preserves balance,
symmetry, and even outdegree at fixed vertices, since the increment
there is $\outdeg_C+\indeg_C=2\outdeg_C$.

\section{Paired walks, sorted layers, and shared axes}\label{sec:layers}
\subsection{Paired decomposition and the exit convention}\label{sec:paired-walks}
At each nonfixed vertex choose a matching from incoming to outgoing
edge copies and force its reversed matching at the mate vertex. At a
fixed vertex, pair outgoing copies by a fixed-point-free involution;
identifying incoming copies with their reversals gives the transition
matching. This is possible by Section~\ref{sec:base}. The resulting cycles
come in distinct reversed pairs. For if reversal fixed one cyclic
edge order, its reflection would fix an edge or exchange two consecutive
edges, yielding a fixed transition. Both are excluded. Fixed tie orders
make this construction deterministic. It is an elementary adaptation
of paired-decomposition principles, not a separate contribution.

Lift each closed walk to windows $[x,e)$ in the periodic text formed
by its cyclic up letters. Up steps increase $e$ and down steps increase
$x$; a turn with $P$ up steps advances by $(P,P)$. A word occurrence
is recorded if a visited window contains it, counted once per starting
residue modulo $P$. The telescoping up-step argument of
\src{\texttt{lem:recorded}} gives total recorded count $m(s)$,
even when $|s|>P$.

Group texts up to translation, reduce to least period $p$, and sort
primitive passages by normalized exits. More explicitly, if original
walk $j$ has turn $d_jp$ and exit function $Z_j$, sort the multiset
\[
\{Z_j(x+rp)-rp: j,\ 0\le r<d_j\}
\]
at each $x$. Its order statistics are layer exits $z_i(x)$.
We index higher layers first. Put
\begin{equation}\label{eq:profiles}
f_i(x)=z_i(x-1),\quad l_i(x)=z_i(x),\quad
z_i(x-1)\le z_i(x),\quad z_i(x-1)\ge x,\quad
z_i(x+p)=z_i(x)+p.
\end{equation}
A layer visits exactly $[x,e)$ for $f_i(x)\le e\le l_i(x)$, then
takes a down step. The source sorting proof preserves the phase-labelled
edge multiset: down steps are determined by exits, and vertical-edge
multiplicities by entry-minus-exit cumulative counts. Each layer has
$p$ up steps and receives an additional budget $p$; the budgets sum
to $W$. Every layer window spells a word of $V$.

The convention for occurrences is crucial. At a specified start $t$,
a layer records $[t,t+h)$ exactly when its \emph{exit}
$z_i(t)\ge t+h$. Thus sorting preserves phase-specific recorded counts,
not merely aggregate counts of word labels. Its entries decide whether
that recorded word is also an actual vertex at start $t$.

\subsection{Text orbits and adjacent pairs}\label{sec:axes}
A nonself text orbit comprises two translation classes exchanged by
$\rho$. Their sorted layers pair rank by rank. Choose the canonical
text to be the greatest forward rotation among \emph{both} orientations.
A basic unit consists of one representative layer and its reflected
mate, with budget $2p$.

For a self-reversed primitive text fix
\begin{equation}\label{eq:reflection}
\mathcal R(x,e)=(c-e,c-x),\qquad
\cm(A(j))=A(c-1-j).
\end{equation}
There are an even number of primitive passages, inherited from distinct
reversed walk pairs. Reflection preserves the order of staircase
profiles: it sends the upper region of a path to the upper region of
the reflected path. The ordered decomposition is uniquely determined
by its exit multisets. Its phase-labelled multiset is invariant under
$\mathcal R$, so every sorted profile is individually fixed under this
reflection, modulo period translation. This does \emph{not} imply even
multiplicity of each profile.

\begin{lemma}[Shared-axis adjacent pairs]\label{lem:axes}
Pair the sorted layers of a self text from the top in adjacent pairs.
Each pair shares both reflection-axis vertices in a primitive turn,
and is therefore already connected. The two profiles may be different.
Removing full pairs preserves all fixed-phase even tails.
\end{lemma}
\begin{proof}
For a nonempty interval $[t,t+h)$ with $2t+h\equiv c\pmod p$,
reflection fixes its alignment. It pairs its recorded occurrences on
the original distinct reversed walks without fixed occurrences.
Consequently their number is even. The exit convention makes this
number exactly
\begin{equation}\label{eq:even-tail}
\#\{i:z_i(t)\ge t+h\},
\end{equation}
which sorting preserves. Since exits are ordered, every even upper
tail consists of whole top-aligned adjacent pairs. Whenever the upper
member visits this axis vertex, the lower member records it too; its
entry is no higher than the upper entry, so it actually visits the same
vertex. The two threshold indicators agree for every pair, and removing
a full pair preserves evenness.

Along one primitive turn, $x+e$ increases by one at every step and by
$2p$ overall. It therefore reaches both classes $c$ and $c+p$ modulo
$2p$, at actual vertices. These are the two fixed-axis vertices, where
the preceding argument applies. If the vertex is empty, its upper
entry is $x$, and the lower entry, lying between $x$ and that entry,
is also $x$; both visit $\eps$. No occurrence count of the empty word
is needed. When $p=1$, every window is an axis window and the turn has
exactly two window visits; the same argument includes this case.
\end{proof}

\subsection{Spliced paths for hosting}\label{sec:splice}
Let $D_1,D_2$ be a self pair and let $U,V$ be its two shared axes in
cyclic order. Form $E$ from $D_1[U,V]$ followed by
$D_2[V,U+(p,p)]$, and form $F$ from the complementary arcs.
These are monotone period-$p$ paths with $p$ up steps each; their
labelled union equals that of $D_1,D_2$. Reflection exchanges the
half-period arcs, so $\mathcal R(E)=F$ at the level of windows and
phase-labelled edges. Coincident profiles use distinct labels.
This is not an assertion about the original base edge-copy involution
after sorting. Later rootedness uses shared vertices, and mirrored
added walks receive new paired edge copies. Keep the original sorted
profiles for all band geometry; use $E,F$ only as host paths.

The exchange-of-halves operation has the antecedent in
\cite[Section 5]{gk}. That operation assumes the two meetings.
Lemma~\ref{lem:axes} is what establishes them for the ordered profiles here.

\section{Geometric connection contracts}\label{sec:operations}
A piece is \emph{rooted} when its underlying undirected support is
connected to $\eps$. All edge multisets are balanced, and all base edges
remain, so touching one vertex of a layer connects its entire support.
Incidental connectivity does not change budget marks.

A \emph{record} is an actual window on a target layer, together with
an alignment in a source text agreeing throughout that window.
The matching-window contract \src{\texttt{lem:matching-window}} says
that the number of blocked recorded occurrences relative to $A$ is
$B_A(s;m)$, and every unblocked occurrence is contained in such an
actual matching window. The nearest-mismatch correspondence proves
the first assertion. For the second, follow the layer from the first
window containing the occurrence until its left mismatch is deleted;
a right mismatch cannot enter first without recording its bracket.

We also use \src{\texttt{lem:agreement}}: distinct aligned texts of
periods $p,q$ agree on fewer than $p+q$ letters, and on fewer than $p$
if both have period $p$. Alignment matters, even within one translation
class. The maximal-rotation contract \src{\texttt{lem:max-rotation}}
says that if $t$ starts the greatest forward rotation of $A$ and every
rotation of a different aligned period-$q$ text $D$ is at most that
word, they cannot agree on all $q$ positions starting at $t$.
These comparisons use encoded-symbol order and include opposite
orientations as well as different phases.

\begin{lemma}[Ordered windows and bands]\label{lem:windows}
For windows $[a,b),[c,d)$ in a common text, spelling words in $V$,
with $a\le c$ and $b\le d$, there is a hierarchical walk of up cost
$d-b$, with a contact of length at most $\max(0,b-c)$.
An ordered list costs its total end increase. If the last window is
a period translate of the first by $K$, it closes at cost $K$.
For any ordered band of period-$p$ layers with top first end $H$
and bottom last end $E$ at $t$, a nonnegative multiple $K$ of $p$
with $H-E\le K$ joins the band at cost at most $K$.
\end{lemma}
\begin{proof}
These are \src{\texttt{lem:ordered-windows}, \texttt{lem:band}}.
If $c\le b$, delete to $[c,b)$ and append to $[c,d)$, using suffixes
and prefixes. If $c>b$, delete to empty, traverse each intervening
letter by an up/down pair, then build the second window. The active
alphabet's individual letters lie in $V$. Thus all intermediate vertices
are legal. For a band, if $E\ge H$ the layers already share $[t,E)$.
Otherwise grow from $[t,E)$ to $[t,H)$, touching each vertical range,
and continue to $[t+K,E+K)$. The end increase is $K$.
\end{proof}

\begin{lemma}[Collective proof in geometric form]\label{lem:collective}
Let any positive number of ordered period-$p$ layers have top first
end $r_H=f_H(a)$ and bottom last end $r_L=l_L(a)$.
Let $R=[a,b)$ be a record and $K$ a nonnegative multiple of $p$.
If
\begin{equation}\label{eq:collective}
r_H\le b+K,\qquad r_H-r_L\le K,\qquad r_L+K\ge b,
\end{equation}
a cost-$K$ closed walk based at $R$ touches all these layers.
\end{lemma}
\begin{proof}
This is the proof of \src{\texttt{lem:collective}}, with its numerical
layer-count restriction removed, not a verbatim use of its statement.
Set $Y=\max(b,r_H)$. The inequalities give $b\le Y\le b+K$ and
$f_i(a)\le Y\le l_i(a)+K=l_i(a+K)$ for every layer.
Choose its first start $x_i\in[a,a+K]$ reaching height $Y$.
Then $[x_i,Y)$ is an actual window. Layer order makes these starts
nondecreasing from top to bottom. Follow $R$, these windows, and $R+K$
using Lemma~\ref{lem:windows}. No step used the number of layers.
\end{proof}

\subsection{The individual trichotomy}\label{sec:individual}
For a period-$p$ layer with $f_i(a)\le b$ and a different-text record
$R=[a,b)$ of target period $q$, the source's Section 5 gives three cases.
If a vertex has length at most $p$, root it by a round trip through
$\eps$. Otherwise, if $l_i(a)+p\ge b$, choose an actual
$[a+p,e')$ with $b\le e'\le b+p$ and follow
\begin{equation}\label{eq:individual}
R,\quad[a+p,e'),\quad R+p.
\end{equation}
This cost-$p$ link has a contact of length at most
$\max(0,b-a-p)\le q$. We use the weak bound $\le q$ in cycle accounting.
If neither case holds, reserve budget $p$ for the actual requester
window $[a,e)$, $e=f_i(a)$, which satisfies
\begin{equation}\label{eq:request}
a+p<e\le l_i(a)<b-p.
\end{equation}
Since $b-a<p+q$ and now $b-a>2p$, necessarily $q>p$.
A request may spend the recipient's budget; a planned link does not.

\begin{lemma}[Single-path host contract]\label{lem:host}
Let a monotone period-$q$ host path have available budget $q$.
For each of a nonempty set of distinct child text groups choose one
actual host record $R_i=[a_i,b_i)$, agreeing with the child text of
period $p_i$. The child texts are distinct up to translation from one
another and from the host text. Every requester at $R_i$ is an actual
child window $[a_i,e)$ satisfying
$a_i+p_i<e\le b_i-p_i$, with its budget $p_i$ reserved.
Then the host and all requesters can be rooted within those budgets.
\end{lemma}
\begin{proof}
This is \src{\texttt{lem:host}}; we recall its constructive content
because it specifies what a spliced host must satisfy. Choose a child
$j$ of minimum period among the requesting groups and order the records
from $R_j$ to its next copy on the host. Translations by multiples of
$q$ transport each entire child coordinate system, preserving its labels
and budget. Put $T=a_j+q$. Then
\[
a_j\le a_i\le T,\qquad b_j\le b_i\le b_j+q.
\]
Successive records obey $b_i-a_l<p_i+p_l$: if they overlap this is
the agreement bound for distinct child texts; if not it is immediate.
The initial host excursion visits all requester windows at $j$, ordered
by end, then $R_j$. At each later record visit the clamped contacts
\begin{equation}\label{eq:clamp}
Q_{i,e}=[\min(a_i+p_i,T),\max(e,b_j)).
\end{equation}
They lie inside $R_i$. Before clamping, their starts and ends are ordered
because $a_i+p_i<e_i\le b_i-p_i<a_l+p_l<e_l$;
clamping preserves order. Each end is less than $T+p_j$, using the
overlap with the next $j$ record. Return to the first requester at $j$
translated by $q$, whose end exceeds $T+p_j$.
Lemma~\ref{lem:windows} gives a cost-$q$ closed excursion and, on the final
transition, a contact of length less than $p_j$.

For a requester at $i\ne j$, write $Q_{i,e}=[s,h)$. Then
$a_i\le s\le a_i+p_i$ and $e\le h\le e+p_i$.
For the last bound, if $h=b_j$, minimum $p_j\le p_i$ gives
$b_j<a_i+p_j+p_i\le a_i+2p_i<e+p_i$.
In the child text the ordered list
$Q_{i,e},[a_i+p_i,e+p_i),Q_{i,e}+p_i$
attaches the requester at cost $p_i$. The requesters at $j$ were touched
directly, so one unspent $p_j$ budget roots the short contact.
This proves the contract, including the case with no later records.
\end{proof}

\section{Routing requests through paired hosts}\label{sec:hosts}
Process primitive-text orbits in increasing least period, breaking ties
by canonical text and stable identifiers. Requests will go strictly
upward in period. All records refer to the fixed base-layer decomposition.
A basic pair reserves its entire $2p$ budget, and processing never splits
that account.

For a nonself recipient pair use its two exchanged sorted layers $E,F$.
For a self pair use the spliced paths of Section~\ref{sec:splice}, each with
budget $q$ and together representing the recipient's $2q$ account.
If requests arrive, choose one original request as an anchor and name
$E$ so that its original record lies there. Keep its whole child-orbit
batch unflipped. Every other child orbit also uses one record for all
its requests. If that record lies only on $F$, reflect the whole batch,
so its record lies on $E$. Distinct child orbits still give distinct
child texts, all of period smaller than $q$.

\subsection{Reclassification before invoking the host lemma}\label{sec:reclassification}
A reflected requester may be right-anchored. At its new record
$R'=[a',b')$, its reflected actual window is $[u,b')\subseteq R'$.
Monotonicity gives $f(a')\le f(u)\le b'$. Reapply
Section~\ref{sec:individual} to that layer and $R'$. Its first case is vacuous
here: the original case-three requester had no vertex of length at
most $p$, and reflection preserves lengths. Its second case gives a
link to $R'$ which will be executed \emph{after} the host is rooted.
Its third case gives exactly the left-anchored inequality in
Lemma~\ref{lem:host}. The unchanged anchor batch still supplies at least one
genuine third-case request.

Apply Lemma~\ref{lem:host} once on the actual monotone path $E$ to the
surviving third-case requests and mirror the entire construction on $F$.
The one-record-per-child, distinct-text, left-anchor, nonempty, and
available-budget hypotheses have all been re-established; this is not
an invocation on a disconnected union of profiles. Choose the
minimum-period child only among these surviving groups. One of its
requester-pair budgets pays the two short root contacts; that pair was
touched directly by the two host excursions. Total cost is at most
$2q$ plus all surviving requester-pair budgets.

Both host paths are rooted. Their union contains all of the original
recipient pair. Now execute the second-case attachments with their own
$2p$ accounts. They are attachments to rooted material, not new outgoing
plans at a noncanonical phase. Thus all original requester units become
rooted. A self child pair is connected at its axes, so touching its
selected member roots both; a nonself pair uses the mirrored construction
for its other member. This establishes the routing without a new
geometric host lemma or separate novelty claim.

\subsection{Two roles of the baseline}\label{sec:baseline}
Choose the highest recipient \emph{pair} as baseline if there are any
recipients. Its two members have already been rooted. The unchanged
anchor gives an actual record $S$ of length $<q+p<2q$ on one member,
where $p<q$ is a child's period. The short-record member and the upper
sorted member used to bound the available region need not coincide.

Root every unspent pair below the baseline with its own $2q$ budget.
If a layer has a vertex of length at most $q$, use that root excursion
and its mirror. Otherwise, writing the short record as $S=[a,h)$,
a lower layer has an actual $[a,e)$ with
$q<e-a\le h-a<2q$. The ordered list
\[
[a,h),\quad[a+q,e+q),\quad[a+q,h+q)
\]
costs $q$ and attaches the lower layer; add its mirror. For a nonself
orbit, if $S$ lies on the partner layer, do this construction in that
partner text group with the same-rank lower member and then reflect.
For a self pair all lower pairs lie below whichever member carries $S$.
The upper member of the baseline pair is used subsequently for band tests.
These two uses of the baseline have distinct hypotheses and budgets.

\section{Upper groups and the matched odd core}\label{sec:groups}
After incoming requests and lower pairs have been handled, the available
upper region has intact pair boundaries. Its budgets are untouched.
We now either root its pairs, reserve outgoing links, or send requests
using one chosen record per child text orbit.

\subsection{Nonself orbits: the ordinary procedure and its mirror}\label{sec:nonself}
Use the canonical representative group. If it has $n$ available layers,
let $t$ start its greatest forward rotation and $H=f_H(t)$.
When $H-t\le np$, an excursion through $[t,H)$ roots them within $np$.
When a baseline exists with $l_{\rm base}(t)\ge H-np$, a band of width
$np$ attaches them. Mirror the chosen operation; the paired cost is
at most $2np$.

In the remaining case $H>t+np$ and the baseline, if present, lies below
$H-np$. Choose the first $k$ with
$\#\{i:l_i(t)\ge H-kp\}\le k$, counting the whole representative group.
It stops by $n$. The preceding failed threshold gives the bottom of
the top $k$ layers last end at least $H-(k-1)p$; for $k=1$ this is
immediate. Translation to $(t-p,t]$ gives the wide band $\le kp$ at
every start. The ordinary $k+1$ rule forces an unblocked occurrence
of $w=A[t:H-kp)$ at a different full-text alignment, and hence a record
$R=[a,b)$ with $a\le t$, $b\ge H-kp$, $|R|<p+q$.
Only canonical-group layers can carry the identical alignment.

Set $K=kp$, $r_H=f_H(a)$, and let $r_L$ be the bottom last end of
that top block. If $r_L+K\ge b$, reserve its $K$ budget for the
collective link. Otherwise use the individual trichotomy on all upper
layers. If the collective link is made, use the trichotomy on the layers
below it. In either case their first ends at $a$ are at most $b$:
when the test fails $r_H\le r_L+K<b$, and below a successful block
$l_i(t)<H-K\le b$. Mirror every root, link, and request with its
paired layer budget. This is precisely the ordinary procedure of
\src{Section 5}, with paired accounting. A collective block and its
mirror form a unit of budget $2kp$; band joining each component costs
at most $2(k-1)p$, leaving $2p$. An individual pair has the same saving
with $k=1$. The canonical orientation does not change with host routing.

\subsection{Self orbits: threshold and odd core}\label{sec:self}
Suppose the available upper region has $2r$ layers, paired from the top.
If $r=0$ there is nothing to do. Take the greatest forward phase $t$
and $H=f_H(t)$. If $H-t\le rp$, an excursion through the top vertex
and its mirror roots all upper layers for at most $2rp$.
If a baseline exists with $l_{\rm base}(t)\ge H-rp$, a width-$rp$
band and its mirror do the same. Otherwise
\begin{equation}\label{eq:self-hard}
H>t+rp,\qquad l_{\rm base}(t)<H-rp\quad\text{if a baseline exists}.
\end{equation}

\begin{lemma}[Odd core and outside record]\label{lem:odd-core}
Choose the first $k\ge1$ such that
\begin{equation}\label{eq:threshold}
\#\{i:l_i(t)\ge H-kp\}\le2k,
\end{equation}
counting the whole self group. Then $k\le r$, and the top $2k-1$
layers, called the core, satisfy
\begin{equation}\label{eq:core-band}
l_{\rm core}(t)\ge H-(k-1)p,\qquad
f_H(x)-l_{\rm core}(x)\le kp\quad(x\in\ZZ).
\end{equation}
For $w=A[t:H-kp)$ there is an actual record $R=[a,b)$ on a different
aligned text of least period $q$, with
\begin{equation}\label{eq:self-record}
a\le t,\quad b\ge H-kp,\quad b-a<p+q.
\end{equation}
\end{lemma}
\begin{proof}
By equation~\eqref{eq:self-hard}, no layer at or below the baseline qualifies
at $k=r$, so the threshold stops by $r$. If $k>1$, failure at $k-1$
means at least $2k-1$ layers qualify there, giving the first bound.
If $k=1$, the core is just the top layer, whose last end is at least
$H$. For $x'\in(t-p,t]$, monotonicity and periodicity give
$f_H(x')\le H$ and $l_{\rm core}(x')\ge l_{\rm core}(t)-p$.
Translation proves the second bound.

The actual top vertex $v=A[t:H)$ is unblocked in its own text, so its
count-rule trigger holds. The self-text rule supplies at least $2k+1$
unblocked occurrences of $w$. At most $2k$ have the identical full-text
alignment, by equation~\eqref{eq:threshold} and the exit convention: each primitive
layer has just one start residue with that alignment. An occurrence
outside it gives an actual matching record by Section~\ref{sec:operations}.
Agreement bounds give equation~\eqref{eq:self-record}; if $q=p$ the stronger
bound $b-a<p$ holds. The alternative alignment may be another phase
of the same group.
\end{proof}

\subsection{Completing a paired unit}\label{sec:core-completion}
Set $K=kp$, $r_H=f_H(a)$, and $r_L=l_{\rm core}(a)$.
Then equation~\eqref{eq:core-band} and equation~\eqref{eq:self-record} give
$r_H-r_L\le K$ and $r_H\le H\le b+K$.
If $r_L+K\ge b$, apply Lemma~\ref{lem:collective} to the $2k-1$ core
layers, at cost $K$, and add its reflection. The final, $2k$th layer
shares the axes with core layer $2k-1$ by Lemma~\ref{lem:axes}.
The whole $2k$-layer unit is therefore connected to its target within
its $2kp$ budget. Alternatively, the sharp core band
in equation~\eqref{eq:core-band} joins the core at cost at most $(k-1)p$;
adding its mirror and using the shared axes internally joins the full
unit for at most $2(k-1)p$, freeing $2p$.
This includes $k=1$: its core is one layer and its mate already meets it.

If the collective test fails, $r_H\le r_L+K<b$, so \emph{all} available
upper pairs, including the top $k$ pairs, enter the individual procedure.
If it succeeds, only pairs below the top $2k$ layers do so; their last
ends at $t$ are $<H-K\le b$, hence their upper first ends at $a$ are
at most $b$. Apply Section~\ref{sec:individual} to each such upper member,
then mirror. Roots and links use at most $2p$; otherwise reserve a
$2p$ request to period $q>p$. A self pair's two members already meet,
so touching the chosen member is sufficient. Every request from this
child orbit uses the \emph{one} record supplied above.

\begin{proposition}[Disjoint paired budget partition]\label{prop:partition}
After processing all text orbits, each basic pair is marked rooted or
belongs to exactly one outgoing unit. A unit targets one actual layer
and its reflected partner, and reserves only its own paired budget.
A unit with budget $2kp$ can be joined internally for at most
$2(k-1)p$. All requests have been fulfilled.
\end{proposition}
\begin{proof}
Incoming requests use disjoint reserved child budgets and the still
available recipient-pair budget. The baseline step spends only remaining
lower-pair budgets. Easy branches spend only available upper budgets;
in the hard branch, the collective unit and individual pairs partition
them. Every request goes to a strictly larger period, so it is fulfilled
later by Section~\ref{sec:hosts}. Only unexecuted links remain outgoing.
The internal saving is the ordinary sharp-band saving for nonself
blocks and Section~\ref{sec:core-completion} for self blocks. Whole-pair removal
preserves fixed-phase parity and pair boundaries. No incidental
intersection is used to reassign or spend another account.
\end{proof}

\section{Opening quotient cycles symmetrically}\label{sec:cycles}
Make one node for each outgoing unit. Its target pair belongs to another
unit or is rooted, so every nonsink node has one outgoing arc.
A nonself unit has two ports, its representative block and its reflected
block; each is internally connectable. A self unit is internally
connected once its core is joined, so its conceptual ports coincide.
A planned link pair joins corresponding source and target ports,
possibly exchanging them. A quotient path therefore lifts to paths
reaching endpoint ports. A twisted quotient cycle need not lift to two
disjoint cycles, and we do not assume that it does.

\subsection{Selfloops and maximal contacts}\label{sec:contacts}
For a quotient selfloop, source and target have equal period $p$ and
different aligned text, even if the target is the opposite orientation.
Thus $|R|<p$. Omit both planned links, join the unit internally, and
use the saved $2p$ to root $R$ and $\rho(R)$ for $2|R|<2p$.
Both ports are rooted.

For a cycle of length at least two, choose a unit with greatest canonical
forward text, breaking ties by stable IDs. Every forward rotation of
\emph{either} orientation of the target text is at most this chosen
source. Let $q$ be the target period.

\begin{lemma}[Paired short contact]\label{lem:cycle-contact}
The selected link pair can be realized within its reserved budget with
one contact of length at most $q$ and its reflection.
\end{lemma}
\begin{proof}
For an individual plan, equation~\eqref{eq:individual} gives the bound $\le q$.
For a collective plan consider its representative block in the nonself
case, or its odd core in the self case. Write $K=kp$, with record
$R=[a,b)$, distinguished start $t$, and $H=f_H(t)$.
If a constituent already has a vertex of length at most $q$, the
collective construction suffices. Otherwise all constituent vertices
have length greater than $q$.

The maximal-rotation contract and the canonical comparison give
$b-t<q$. Put $x=\min(t,a+K)$. If $x=t$ this already gives $b-x<q$;
otherwise
$b-x=b-a-K<p+q-K\le q$. Thus $[x,b)$ is a short suffix of $R$.
Every constituent first end at $x$ exceeds $b$, since its vertex length
exceeds $q>b-x$, and is at most $H$. With $F=f_H(x)$ we have
$b<F\le H\le b+K$. Follow the ordered list
\begin{equation}\label{eq:cycle-list}
[a,b),\quad[x,b),\quad[x,F),\quad[a+K,b+K).
\end{equation}
Its cost is $K$; the growth at $x$ touches all constituents, and the
short suffix is the desired contact. Its prescribed windows are actual
vertices or substrings of them. This is the geometric proof of
\src{\texttt{lem:cycle-contact}}, now applied to the odd core where
appropriate, not an unqualified use of its original block statement.
The final self-layer meets its core partner. Reflect the constructed
walk to obtain the second contact within the paired budget.
\end{proof}

\begin{theorem}[Paired connection budget]\label{thm:connection}
One can add paired closed walks of total up cost at most $W$ to $G$
so that the resulting graph $D$ is connected to $\eps$. It is balanced,
$\tau$-invariant, and has even outdegree at every fixed vertex.
Its total up cost $U$ is at most $2W$.
\end{theorem}
\begin{proof}
For a cycle of length at least two retain the selected link pair of
Lemma~\ref{lem:cycle-contact}, but omit the \emph{target unit's} outgoing
pair. These are different links. Join that target internally, saving
$2q$, and use the saving to root both short contacts. The quotient
cycle with the omitted edge removed is a path ending at the target.
Every lifted component meets one of the retained contact pieces, which
touch both target ports. Rootedness propagates along the path to every
source port, whether the lift was twisted or untwisted. This also
covers mixed self and nonself cycles.

Different quotient cycles have disjoint accounts. Execute all other
links in increasing distance from rooted sinks or opened cycles.
Each uses its source unit's reserved account. Together with the earlier
rooting and request operations, this roots every layer. The accounts
partition the original layer budgets, whose sum is $W$.

Every addition was a closed walk and its reversed partner, with fresh
paired edge copies. Hence balance, symmetry, and fixed-vertex parity
persist by Section~\ref{sec:base}. All base edges remain, so required vertices
remain used, and total up cost is at most $W+W=2W$.
\end{proof}

\section{Half-Euler extraction and polynomial bit complexity}\label{sec:extraction}
\begin{lemma}[Half-Euler walk]\label{lem:half-euler}
Let a finite nonempty balanced directed multigraph have a reversing
involution on vertices and edge copies, no fixed edge, connected used
support, and a fixed root in that used support. If every fixed vertex has even outdegree,
there is a directed closed walk from the root using exactly one edge
from each involution pair. It can be constructed in polynomial time.
\end{lemma}
\begin{proof}
Quotient vertices and edges by the involution, retaining both incidences
of each edge, including loops. At a nonfixed vertex orbit choose a
representative $v$. Label an incidence positive if its compatible lift
is outgoing at $v$, and negative if incoming there. The numbers of
positive and negative incidences equal $\outdeg(v)$ and $\indeg(v)$,
so they agree. At a fixed vertex the quotient degree equals its directed
outdegree and is even.

Pair opposite signs at each nonfixed orbit and pair arbitrary incidences
at fixed vertices. The resulting transition system decomposes the
quotient edges into closed alternating trails. These lift consistently:
at a nonfixed orbit an arrival through a positive incidence puts the
lift at the mate of $v$, where departure requires a negative incidence;
the other case is reversed. At a fixed vertex either state is available.
The final/initial paired incidences give closure. If several trails
remain, connectedness supplies two sharing a vertex. Cross-switch two
incidence pairs there, aligning opposite signs at a nonfixed orbit,
to merge the trails. This preserves all local conditions. Repeating
produces a quotient Euler trail, which lifts from the fixed root and
closes there, selecting one copy from each orbit.

This is the signed-incidence Euler construction underlying the
antecedents in Section~\ref{sec:antecedents}, with fixed vertices handled
explicitly. The general lemma allows quotient loops. Our hierarchical
graph has none, because an edge changes word length and $\rho$
preserves it.
\end{proof}

\begin{proof}[Proof of Theorem~\ref{thm:main}]
The empty cases were handled before the construction. Otherwise apply
Lemma~\ref{lem:half-euler} to $D$ from Theorem~\ref{thm:connection}, rooted at
$\eps$. Since $D$ has $U$ up and $U$ down edges and $\tau$ exchanges
them, it has $U$ edge orbits. The selected closed walk has $U$ edges;
half are up, since word length returns to zero. It visits a representative
of every used vertex orbit, and thus one orientation of every required
word. Append the letter of each up step to the output. At every moment
the current vertex is a suffix of the output so far: up appends a letter
to both, and down deletes the vertex's first letter. This is the suffix
invariant of \src{Section 7}. Therefore the word is feasible and
\begin{equation}\label{eq:final}
|T|=U/2\le W\le2\OPT(\mathcal S).
\end{equation}

For completeness we detail the finite realization and bit costs.
All finite choices use encoded-symbol order and stable IDs, including
orbits of equal least period, tied canonical maxima, transitions,
records, host anchors, threshold choices, and trail splices.
Unused alphabet symbols may be ignored after reading and checking the
explicit alphabet and involution; keep every used symbol and its mate.

There are $O(L^2)$ substring vertices of length at most $L$, and all
counts are at most $2L$. Enumerating $v,p,k$ as in Section~\ref{sec:blocking}
gives a polynomial list (the coarse bound $O(L^4)$ suffices).
Evaluate blocking by enumerating every bracket word and position.
Self-reversal is tested by comparing a primitive period word with
rotations of its reverse complement. These are polynomial operations
on explicit strings and polynomial-bit integers.

The base graph has $2W\le4L$ edge copies. Paired transitions and walk
lifts are finite scans. Least periods, normalized passages, and their
exit order statistics use explicit comparisons and sorting. The total
number of primitive phases is $W$, so storing one period of each
$z_i$ and its text is polynomial. Both axis crossings and the host
splices are scans through these stored paths.

Equality or order comparison of two periodic texts needs at most the
sum of their periods, by the agreement bound. Normalize record starts
to a turn, translating child coordinates together with the record.
All actual windows have length at most $L$, and the resulting integers
have polynomial bit length. There are $2W$ actual window visits per
turn across all layers. To find a different-alignment record for $w$,
enumerate them and each placement of $w$ within them, at most $2WL$
candidates. Align the placement to its prescribed occurrence interval;
reject identical full texts or windows not matching the source throughout.
The count and matching-window arguments prove that a candidate remains.

Threshold searches, the individual tests, batch reclassification, host
record order, and functional-graph cycle detection all operate on
polynomially many stored objects. Collective intersections scan at most
$K\le W$ start coordinates. Ordered-window paths are materialized by
prefixes, suffixes, and single letters. The paired budget proof bounds
the final graph by $2U\le4W\le8L$ edge copies, including all additions.
The quotient traversal and output take polynomial time in this explicit
graph. Comparisons and output of variable-width symbol labels are
charged their encoded bit lengths; no unbounded-label unit-cost operation,
random choice, or infinite-text oracle is used. This establishes the
claimed deterministic polynomial bit complexity at the full scope.
\end{proof}

\appendix
\section{Scope of formal support and validation}\label{app:formal}
A companion Lean package, \texttt{RCS128 formal v1}, contains sixteen
support theorems. It does not formalize Theorem~\ref{thm:main}. Its independent
formal review concerns these exact finite or scalar statements, with
standard axioms only. The correspondence is as follows; theorem names
are in namespace \texttt{RCS128}.

\begin{longtable}{@{}>{\raggedright\arraybackslash}p{.44\linewidth}>{\raggedright\arraybackslash}p{.52\linewidth}@{}}
\toprule
Declarations & Mathematical interface and limitation\\\midrule
\endhead
\path{even_card_of_reflection}, \path{even_phase_card}
& Even cardinality under a supplied free involution, optionally preserving
a phase predicate. The actual tagged occurrences and their reflection
are not constructed in Lean.\\[4pt]
\path{filter_eq_range_at_boundary}, \path{adjacent_pair_of_even_tail},
\path{adjacent_exit_threshold}, \path{reflection_sorted_pair}
& Even upper tails force adjacent top-aligned indicators to agree.
The exact equality between phase counts and exit counts, sortedness,
and phase preservation are hypotheses.\\[4pt]
\path{even_card_erase_pair}, \path{even_phase_after_pair_removal},
\path{lower_pair_visits_axis}
& Removal preserves parity when both distinct elements belong to the
underlying set and agree on the phase; entry order upgrades a recorded
threshold to a visit. Axis existence and profile geometry are external.\\[4pt]
\path{odd_core_from_previous_failure}, \path{first_threshold_core},
\path{exists_record_outside_alignment}, \path{odd_core_budget}
& Minimal-threshold arithmetic, the $2k+1$ versus $2k$ pigeonhole
step, and $2(k-1)p+2p=2kp$. Actual records, bands, and the $k=1$
geometric case are external.\\[4pt]
\path{shifted_copy_fits}, \path{fixed_deficit_surplus}
& Offset fitting and an even positive deficit giving at least $2k+1$.
The reflected bracket construction and complete domination induction
are external.\\[4pt]
\path{final_length_bound}
& Scalar composition of $2|T|=U$, $U\le2W$, and $W\le2\OPT$.
All three substantive premises are supplied outside Lean.\\\bottomrule
\end{longtable}

\subsection{What the support does not establish}\label{app:formal-limits}
The frozen package does not formalize words, reverse complements,
blocking, hierarchical graphs, sorted staircase geometry, spliced hosts,
quotient cycles, half-Euler extraction, or polynomial bit complexity.
The unblocked and nonfixed-bracket arithmetic $(k+1)+k=2k+1$ is not
part of its frozen count support. Removing pairs, re-indexing the sorted
list, and pairing the remainder is not formalized as a composed layer
operation. In the removal statements, ``absent pairs'' means absent
from the \emph{phase}; both distinct elements must still belong to the
underlying set. The extra $1<k$ assumption in
\path{odd_core_from_previous_failure} is harmless but unnecessary;
the separate first-threshold theorem includes $k=1$.
Most declarations are elementary finite or scalar lemmas. None should
be described as a machine proof of the geometric construction.

The seven frozen author control declarations are mostly numerical or
finite-set witnesses, not formal hypothesis-sensitivity tests of those
theorems. The independent review added fifty-nine declarations that
supply the sensitivity evidence and positive controls. We attribute
that evidence to the reviewer, not to the original seven controls.

The recorded review rebuilt the support and reviewer sources with Lean
4.34.1 and inspected printed types and axioms. It performed module-only
kernel replay. Pinned dependency \texttt{.olean} files were trusted
read-only; no fresh replay or rebuild of all mathlib dependencies was
performed. Compilation alone does not certify statement identity:
a weakened factor-four statement passes the original harness.
Identity therefore also rests on external manifest hashes and the
reviewer's type audit. The original harness has declaration/token-scan
limitations, closed for these frozen sources by that audit and the
module replay, not by a universal claim about the harness.

\subsection{Diagnostics and author-side provenance}\label{app:diagnostics}
Finite count, axis, splice, and graph checks are diagnostics, not proofs.
An independent proof reviewer also wrote an end-to-end executable model
and tested hosts, cores, cycles, and output. Actual multi-child flipped
batches occurred zero times in the recorded model runs; the routing
argument in Section~\ref{sec:reclassification} rests on its mathematical proof,
not on an assertion that testing covered that branch. Per-request
shadow checks do not change that limitation. The original frozen
candidate, its exact scope, and all failed intermediate constructions
are preserved outside this manuscript; no historical novelty statement
is silently updated there.

\section{Attribution, reading limits, and rights}\label{app:attribution}
This is a source-relative article by Myeongjin Kang, with no affiliation
stated. AI assistance is part of its author-side provenance. The parent
GPT-6 Astra supplied substantial author-side work on the reflection
count argument, odd-core rule, spliced-host routing, cycle accounting,
baseline distinction, empty-vertex degree correction, and scoped Lean
implementation. Author-side agents developed the construction and
checks; GPT-6 Astra prepared this manuscript. Separate Opus 5.5 agents
performed the proof, formal-support, and contribution reviews. These
were independent agent roles, not human refereeing or author-side
checks relabelled as independent review. Those review results are not
publication, acceptance, or worldwide-priority certificates.

The pinned source's sections 1--7 and operative proofs were read for
this manuscript. Primary comparison is bounded: the WABI formulation,
lemmas and algorithm sections, the CPM introduction and Lemma 18,
Goldberg--Karzanov's specified decomposition, parity and splice passages,
and the cited bidirected and rotation passages were inspected.
Fici et al.'s arXiv v1 extraction was read throughout. The accompanying
reading ledger records exact ranges and distinguishes this author's
reading from the parent's broader reading. The Jiang--Li--Du 1992
body and Breslauer--Jiang--Jiang 1997 journal body remain unread under
the explicit citation choices in Section~\ref{sec:antecedents}. We make no
exhaustive literature or worldwide-priority claim.

Original prose is licensed CC BY 4.0 and original accompanying code MIT.
Private third-party source bodies are not included in the public
manuscript payload. Citations identify sources without asserting a
redistribution license for their downloaded bodies.

\begin{thebibliography}{99}
\bibitem{ordinary} OpenAI.
\emph{A Polynomial-Time 2-Approximation for Shortest Common Superstring}.
September 24, 2026. Pinned source, revision
\texttt{fd4aeeb2ee4fc729c18d98444fed42fd0529eeeb}.
\href{https://github.com/openai/math/tree/fd4aeeb2ee4fc729c18d98444fed42fd0529eeeb/preprints/A-Polynomial-Time-2-Approximation-for-Shortest-Common-Superstring-September-24-2026}{Pinned manuscript directory}.

\bibitem{yswabi} Ryosuke Yamano and Tetsuo Shibuya.
Improved approximation algorithms and hardness results for shortest
common superstring with reverse complements.
In \emph{WABI 2026}, LIPIcs 390, 2:1--2:13, 2026.
\url{https://doi.org/10.4230/LIPIcs.WABI.2026.2}.

\bibitem{yscpm} Ryosuke Yamano and Tetsuo Shibuya.
Improved approximation ratios for the shortest common superstring
problem with reverse complements.
In \emph{CPM 2026}, LIPIcs 369, 15:1--15:11, 2026.
\url{https://doi.org/10.4230/LIPIcs.CPM.2026.15}.

\bibitem{fici} Gabriele Fici, Tomasz Kociumaka, Jakub Radoszewski,
Wojciech Rytter, and Tomasz Wale\'n.
\emph{On the Greedy Algorithm for the Shortest Common Superstring
Problem with Reversals}. arXiv:1511.08431v1, 2015.
\url{https://arxiv.org/abs/1511.08431v1}.

\bibitem{jld} Tao Jiang, Ming Li, and Ding-Zhu Du.
A note on shortest superstrings with flipping.
\emph{Information Processing Letters} 44(4):195--199, 1992.
\url{https://doi.org/10.1016/0020-0190(92)90084-9}.
Body unavailable in the bounded comparison; historical discussion is
through \cite{fici,yscpm,yswabi} as explained in Section~\ref{sec:antecedents}.

\bibitem{gk} Andrew V. Goldberg and Alexander V. Karzanov.
\emph{Maximum Skew-Symmetric Flows and Matchings}.
arXiv:math/0304290v2, January 2004 (manuscript dated December 2003).
\url{https://arxiv.org/abs/math/0304290v2}.

\bibitem{medvedev} Paul Medvedev, Konstantinos Georgiou, Gene Myers,
and Michael Brudno. \emph{Computability of Models for Sequence Assembly}.
2007. Author-hosted WABI paper.
\url{https://medvedevgroup.com/papers/wabi07.pdf}.

\bibitem{kundeti} Vamsi Kundeti, Sanguthevar Rajasekaran, and Hieu Dinh.
\emph{An Efficient Algorithm for Chinese Postman Walk on Bi-directed
de Bruijn Graphs}. arXiv:1006.4828v1, 2010.
\url{https://arxiv.org/abs/1006.4828v1}.

\bibitem{bjjjournal} Dany Breslauer, Tao Jiang, and Zhigen Jiang.
Rotations of periodic strings and short superstrings.
\emph{Journal of Algorithms} 24(2):340--353, 1997.
\url{https://doi.org/10.1006/jagm.1997.0861}.
Archival reference; journal body not read in this comparison.

\bibitem{bjjreport} Dany Breslauer, Tao Jiang, and Zhigen Jiang.
\emph{Rotation of Periodic Strings and Short Superstrings}.
BRICS Report RS-96-21, June 1996. Operational comparison version.
\url{https://www.brics.dk/RS/96/21/BRICS-RS-96-21.pdf}.

\bibitem{mucha} Marcin Mucha. Lyndon words and short superstrings.
In \emph{SODA 2013}, pages 958--972, 2013.
\url{https://doi.org/10.1137/1.9781611973105.69}.

\bibitem{collapsing} Alexander Golovnev, Alexander S. Kulikov,
Alexander Logunov, Ivan Mihajlin, and Maksim Nikolaev.
Collapsing superstring conjecture.
In \emph{APPROX/RANDOM 2019}, LIPIcs 145, article 26, 2019.
\url{https://doi.org/10.4230/LIPIcs.APPROX-RANDOM.2019.26}.
Operational comparison: revised arXiv:1809.08669v3, 2020,
\url{https://arxiv.org/abs/1809.08669v3}.
\end{thebibliography}

\end{document}
